\section{A Welder schedule over the tile DSL: conv2d $\rightarrow$ batch norm $\rightarrow$ SiLU}
\label{sec:welder}

This chapter is a design: how Welder (Shi et al., OSDI'23) would schedule a
three-node graph if its \textsc{Propagate} were the DSL's \code{resolve_inputs}. It
follows the paper's algorithms and uses \code{../Welder} where the implementation
departs from the paper in ways that matter. Where teenygrad today cannot do a step,
the chapter says what is missing.

\subsection{Welder in two pages}

\textbf{The model (\S3.1).} A DNN operator is executed as many homogeneous
\emph{operator-tiles}, each computing one output tile from the input tiles it
depends on. Two adjacent operator-tiles can be \emph{connected} through a
\emph{reuse-tile} held at some memory level: registers (L0), shared memory (L1), or
global memory (L2, meaning not connected). The consumer then reads the producer's
tile directly instead of a materialised tensor. A graph of connected operator-tiles
is a \emph{tile-graph}. Because of the dependencies, choosing one output tile
determines every tile in a connected group. \textsc{Propagate} infers them backwards
from the output, by analysing each operator's tensor expression.

\textbf{The interface (Fig.~6).}
\begin{lstlisting}
void       SetConnect(Edge *edge, MemLevel level);               // pick the reuse level of one edge
TileConfig Propagate(TileGraph g, Map<Axis, Dim> config);        // output tile -> all tiles
size_t     MemFootprint(TileGraph g);                            // peak bytes at the level (best-fit)
size_t     MemTraffic(TileGraph g);                              // bytes crossing into the lower level
\end{lstlisting}

\textbf{The insight that makes the search tractable.} Traffic between two levels is
determined by the tile configuration \emph{at that level alone} (``inter-layer
independence''). Each sub-graph at each level can therefore be optimised separately,
with only the constraint that a lower level's tile is larger.

\textbf{The scheduler (Fig.~7).}
\begin{lstlisting}
GraphConnecting(g, d):                          // outer: per edge, per level
  for node in TopologySort(g): for edge in node.out_edges(): for level in d.MemLevels():
      SetConnect(edge, level); s = ExtractSubgraph(node, 0)
      configs = SubGraphTiling(s, 0, shapes)
      if (t = Min(d.Profile(configs))) < best: best = t, best_level = level
    SetConnect(edge, best_level)
SubGraphTiling(g, level, c):                    // inner: tile search at one level
  for subtile in EnumerateSubtiles(g, c):       // Roller-style expansion from small tiles
      config = Propagate(g, subtile)
      if MemFootprint(g) > level.capacity: continue
      configs.push(config, priority = MemTraffic(g))
  for config in TopK(configs): recurse on ExtractSubgraph(node, level + 1) for each sub-graph
\end{lstlisting}

The result is a \emph{hierarchical tile-graph}. \textsc{ExecuteGraph} (Fig.~8) turns
it into code: \code{Allocate} the footprint, \code{LoadTiles}, \code{ComputeTile} each
node (or recurse one level up), then \code{StoreTiles}. \S4.1 adds hardware alignment
to the cost:
\begin{itemize}
\item a penalty on traffic for uncoalesced access;
\item a penalty proportional to idle cores when a tile is so large that there are too
  few tiles (the ``wave'' count);
\item an infinite penalty above capacity;
\item a common thread-block size taken as the GCD of the operators' tile counts.
\end{itemize}

\textbf{Where \code{../Welder} departs from the paper} (\code{ANALYSIS.md} \S7). These
departures matter for what follows.
\begin{itemize}
\item \textbf{Propagation is bound analysis over TVM tensor expressions}
  (\code{shape_inference/te.py}). There is no convolution special case: the halo
  $(T-1)s+k$ falls out of \code{ConstIntBound} on the index expressions
  (\code{SHAPES.md} \S3). Regions are clamped to the producer's real shape.
\item \textbf{Two levels, not a recursion.} Global $\to$ shared $\to$ register is
  hard-wired, and \code{SubGraphTiling}'s recursion and \code{ExtractSubgraph} are
  absent.
\item \textbf{Graph connecting is greedy and measured.}
  \code{Engine._build_fusion_group} grows a group by BFS over output edges, capped at
  10 nodes per step. It accepts a growth only if the \emph{profiled} latency gain is
  positive. The \code{MemTraffic} model is used only inside the tile search.
\item \textbf{Register-level connection is narrow.} It is tried only for ``one reduce
  node followed by one element-wise node'', by splicing the TE graphs
  (\code{compute_inline}).
\item \textbf{The front end pre-digests the graph.} For our graph (\code{SHAPES.md}
  \S1), \code{SimplifyInference} splits BN into a \code{[C]} scale/shift node and an
  element-wise apply. \code{WelderFuseOps} folds conv and its activation into one
  node before scheduling starts, and \code{WelderConvImplicitGemm} turns the conv
  into an implicit GEMM with TensorCore annotations.
\item \textbf{\code{rstep} is a first-class tile dimension.} It is the chunk of each
  reduction axis staged per step, and it controls a conv's halo and shared-memory
  footprint (\code{SHAPES.md} \S4).
\end{itemize}

\subsection{Mapping Welder's primitives onto the DSL}

{\small
\begin{xltabular}{\linewidth}{>{\raggedright\arraybackslash}p{3.0cm}>{\raggedright\arraybackslash}p{3.9cm}>{\raggedright\arraybackslash}X}
\toprule
\textbf{Welder} & \textbf{\code{../Welder}} & \textbf{Over the teenygrad DSL} \\
\midrule
\endhead
Propagate (per node) & \code{IRNode.propogate_inputs}, TE bound analysis & \code{resolve_inputs(spec, tile, consts)}. Name matching plus the receptive field. Implemented. \\
Propagate (graph) & reverse-topo loops in \code{policy/default.py} & A reverse-topo walk passing \code{TileShape}s by dim position, with full-extent materialisation, clamping and re-boxing (\S\ref{sec:edges}). Not implemented. \\
\code{Map<Axis, Dim>} config & output tile + \code{rstep} & Values for the anchor's \code{block_const}s (\code{BLOCK_OW}, \code{BLOCK_N}, \ldots). There is no \code{rstep}: \code{#[tile_loop]} walks reductions one element at a time. \\
EnumerateSubtiles & \code{DFS_smem_tile} from \code{get_base_tile} & Powers of two for each \emph{bound} axis of the anchor's output. Unbound axes are not searchable, because changing them means a different kernel. \\
MemTraffic & boundary reads/writes $\times$ tile count, coalescing factor & $\sum$ (resolved boundary-input bytes + output bytes) $\times$ \#tiles. Needs a graph-level \code{ConstLookup} (\issue{1nr.30}) and the fixes in \S\ref{sec:conv-overestimate}. \\
MemFootprint & best-fit over cached tensors & Bytes of the live tiles in one program, checked against the per-program budget. \\
SetConnect levels & global / shared / register & Effectively two: \emph{global} (separate launches) and \emph{in-program} (fused, so the intermediate is a \code{Tile} value). A Triton-style program does not choose registers vs.\ shared memory explicitly; the compiler does. \\
Block alignment & \code{check_tile_shape_isvalid}, GCD block size & \code{GridSpec} rider check: every fused node's grid axes must be a name-matching subset of the anchor's (\issue{1nr.17}). \\
blockIdx remapping & \code{_assign_block_order} & \code{#[tile_grid(order = ..)]}. With \code{swizzled}, no rider is possible. \\
ComputeTile / codegen & per-node \code{__device__} fns sharing a shared-memory arena & \code{FusionCore} splicing. For a generated loop, the \code{finish} slot (\S\ref{sec:welder-infer}). \\
\bottomrule
\end{xltabular}
}

\subsection{The graph}

The test in \code{anduin/tile_graph/mod.rs} builds the graph from real layers:
\code{Conv2d(3->8, 3x3, stride 1, pad 1)}, \code{BatchNorm2d(8)},
\code{Silu}, over a \code{[10, 3, 32, 32]} input. Lowered in inference mode, it
becomes:

\begin{center}
\small
\begin{tabular}{llll}
\toprule
node & kernel & block & spec rank / bindings \\
\midrule
conv & \code{conv2d_forward} & \code{BLOCK_OW = 16} & 4 / \code{OW}; \code{x}: windowed \code{H} (block 1), \code{W} \\
bn & \code{batch_norm_2d_nchw_forward_inference} & \code{BLOCK_HW = 128} & 4 / \code{HW} over dims [2,3] \\
silu & \code{silu_forward} & \code{BLOCK_SIZE = 1024} & 4 / one flat binding over [0..4) \\
\bottomrule
\end{tabular}
\end{center}

Every intermediate edge carries \code{[10, 8, 32, 32]} f32 tensors (\code{|y|} = 327{,}680 B).
\S\ref{sec:welder-search} also uses a ResNet-stage shape, \code{[32, 64, 56, 56]}
with 64 output channels, where \code{|y|} = 25.7~MB.

\begin{figure}[h]
\centering
\begin{tikzpicture}[
  op/.style={draw=boxline,fill=bg,rounded corners=3pt,minimum width=2.1cm,minimum height=0.85cm,font=\small},
  io/.style={font=\small\itshape,text=cm},
  >={Stealth[length=2mm]}]
\node[io] (in) {x};
\node[op,right=0.9cm of in] (conv) {conv2d};
\node[op,right=2.0cm of conv] (bn) {BN (affine)};
\node[op,right=2.0cm of bn] (silu) {SiLU};
\node[io,right=0.9cm of silu] (out) {y};
\draw[->] (in) -- node[above=1pt,font=\scriptsize]{global} (conv);
\draw[->,very thick,attr] (conv) -- node[above,font=\scriptsize,text=attr]{in-program} (bn);
\draw[->,very thick,attr] (bn) -- node[above,font=\scriptsize,text=attr]{in-program} (silu);
\draw[->] (silu) -- node[above,font=\scriptsize]{global} (out);
\begin{scope}[on background layer]
\node[draw=attr,dashed,rounded corners=4pt,fit=(conv)(bn)(silu),inner sep=5pt,label={[font=\scriptsize,text=attr]below:inference: one virtual node, one kernel (= \texttt{conv2d\_bn\_silu\_forward})}] {};
\end{scope}

\node[io,below=2.4cm of in] (in2) {x};
\node[op,right=0.9cm of in2] (conv2) {conv2d};
\node[op,right=2.2cm of conv2,yshift=0.75cm] (st) {BN stats};
\node[op,right=2.2cm of conv2,yshift=-0.75cm] (nm) {BN normalize};
\node[op,right=2.0cm of nm] (silu2) {SiLU};
\node[io,right=0.9cm of silu2] (out2) {y};
\draw[->] (in2) -- node[above,font=\scriptsize]{global} (conv2);
\draw[->] (conv2) -- node[above,sloped,font=\scriptsize]{global} (st);
\draw[->] (conv2) -- node[below,sloped,font=\scriptsize]{global} (nm);
\draw[->] (st) -- node[right,font=\scriptsize]{global, \texttt{[C]}} (nm);
\draw[->,very thick,attr] (nm) -- node[above,font=\scriptsize,text=attr]{in-program} (silu2);
\draw[->] (silu2) -- (out2);
\begin{scope}[on background layer]
\node[draw=attr,dashed,rounded corners=4pt,fit=(nm)(silu2),inner sep=5pt,label={[font=\scriptsize,text=attr]below:training: three kernels; only normalize+SiLU fuse}] {};
\end{scope}
\end{tikzpicture}
\caption{The schedules derived in \S\ref{sec:welder-infer} and \S\ref{sec:welder-train}.}
\label{fig:schedules}
\end{figure}

\subsection{Inference, step by step}
\label{sec:welder-infer}

\textbf{Step 0: normalise the graph.} Fold BN's four per-channel parameters into
\code{scale = gamma / sqrt(var + eps)} and \code{shift = beta - scale * mean}, which
is Welder's \code{SimplifyInference}. Teenygrad already has the helper,
\code{nn::fused::prefold_bn_affine}. The \code{[C]} fold is a tiny separate node. BN
becomes a pure per-channel affine map, element-wise with two broadcast operands.

\textbf{Step 1: choose the anchor and enumerate.} \code{GraphConnecting} visits
\code{conv} first and tries connecting \code{conv->bn} in-program. The extracted
sub-graph \{conv, bn\} has \code{bn} as its output node, so tiles are enumerated on
\code{bn}'s output. Here the DSL raises its first difficulty: \code{bn}'s only
searchable axis is \code{BLOCK_HW}, while the conv's is \code{BLOCK_OW}. A fused
group has one grid, so the enumerator must pick the axis set of the node whose
\emph{body} is least flexible. That is the conv, whose decode and loop are fixed by
its declaration. The element-wise nodes adapt to any tile. So the anchor for
enumeration is the conv, and the candidates are \code{BLOCK_OW} $\in\{8, 16, 32\}$.
They describe the output tile $[1, 1, 1, b]$ of every node in the group.

\textbf{Step 2: propagate.} The seed is the group output's tile $[1,1,1,b]$, with
$b=16$. Each row is one \code{resolve_inputs} call, and each arrow is an edge
crossing:

\begin{center}
\small
\begin{tabular}{lll}
\toprule
& \textbf{with today's specs} & \textbf{with gridded axes as block 1} \\
\midrule
silu $y$ (seed) & \code{[1, 1, 1, 16]} & \code{[1, 1, 1, 16]} \\
silu $x$ & \code{[1, 1, 1, 16]} (flat, product) & \code{[1, 1, 1, 16]} \\
$\to$ bn $y$ & \code{[1, 1, 1, 16]} & \code{[1, 1, 1, 16]} \\
bn $x$ & \code{[None, None, 1, 16]} & \code{[1, 1, 1, 16]} \\
bn scale/shift & \code{[None]} (all $C$) & \code{[1]} \\
$\to$ conv $y$ & \code{[10, 8, 1, 16]} (full extent materialised) & \code{[1, 1, 1, 16]} \\
conv $x$ & \code{[None, None, None, 18]} & \code{[1, None, 3, 18]} \\
conv $w$ & --- (not in spec) & \code{[1, 3, 3, 3]} (if declared) \\
\bottomrule
\end{tabular}
\end{center}

The right-hand column is what the program actually touches. Two observations:
\begin{itemize}
\item \textbf{The left column's \code{conv y = [10, 8, 1, 16]} is harmless on its
  own.} The conv spec discards it. But a model that charged footprint per edge would
  stage a 10$\times$8$\times$16 slab for the bn$\to$conv reuse-tile when one strip is
  live.
\item \textbf{\code{bn}'s \code{BLOCK_HW} is not used at all inside the group.} The
  rider adopts the anchor's tile. \code{BLOCK_HW = 128} only describes \code{bn} as a
  standalone kernel. That is why the edge crossing must be by dim position, and why
  \code{span} was needed: without it the rank check fails at the \code{bn} row and
  propagation stops.
\end{itemize}

\textbf{Step 3: cost.} The per-program footprint is \code{x}'s region plus the
weights. BN and SiLU add two scalars and no staged tiles. Traffic is the per-tile
boundary bytes times the tile count. Computed with the corrected regions, for the
test shape:

\begin{center}
\small
\begin{tabular}{rrrrr}
\toprule
\code{BLOCK_OW} & tiles & footprint / program & traffic, fused & traffic, three kernels \\
\midrule
8  & 10{,}240 & 468 B   & 5.20 MB & 6.43 MB \\
16 & 5{,}120  & 756 B   & 4.24 MB & 5.51 MB \\
32 & 2{,}560  & 1{,}332 B & \textbf{3.76 MB} & 5.05 MB \\
\bottomrule
\end{tabular}
\end{center}

Every candidate fits, so \textsc{TopK} keeps \code{BLOCK_OW = 32}. On a 32-wide
output that is one tile per row. Fusion saves $4|y|$ ($= 1.31$~MB: BN's read and
write and SiLU's read and write of the intermediate), less the per-tile re-reads of
\code{scale}/\code{shift}. The net saving is 1.23--1.29~MB across the rows.

\textbf{Step 4: connect.} \code{conv->bn} in-program beats global, and so does
\code{bn->silu} on the next visit. The group becomes one virtual node. The rider
check must hold for each edge: \code{bn}'s grid \code{[B (Y), C, HW]} against the
anchor's \code{[B, C_OUT, OH, OW]}. It fails as literally written, for two reasons:
\begin{itemize}
\item \code{C} vs \code{C_OUT}: same dim, different name.
\item \code{HW} vs \code{(OH, OW)}: one flattened axis against two.
\end{itemize}
Since the walk crosses edges by position, the check should too. Dim~1 is dim~1
whatever each kernel calls it. A flattened rider axis is compatible with the anchor
when the anchor's tile is a valid box of it (\S\ref{sec:edges}). SiLU's single flat
axis is a subset of anything. That relaxation (positional rider matching) is the
first concrete change this graph asks of \issue{1nr.17}.

\textbf{Step 5: generate.} The fused program is the conv's prelude and loop, with
BN and SiLU applied to the carry before the store. That is
\code{conv2d_bn_silu_forward}'s \code{finish}, word for word (\S\ref{sec:prelude}).
This suggests the answer to the question \issue{1cli} left open:

\begin{note}[Proposal: a generated-loop kernel's fusion core is its carry]
For a \code{#[tile_loop(generate)]} kernel, the splice point is not ``the body minus
the store'' but ``the carry after \code{finish}''. Fusing a consumer means appending
its core to \code{finish}, with the consumer's \code{input_ident} bound to the carry:
\begin{itemize}
\item \textbf{SiLU} already has a \code{FusionCore}. Its \code{x.tensor} becomes
  \code{acc}.
\item \textbf{BN} does not have one (five inputs). Its broadcast operands would
  become rider loads keyed by the anchor's decoded index. In the hand-written
  kernel that is the \code{arange(0, 1) + tile_c_out} load in \code{finish}.
\end{itemize}
A \code{FusionCore} variant whose extra inputs are broadcast operands, each with the
anchor axis it is indexed by, would make BN spliceable. Since the macro already knows
which inputs are broadcast operands, it could emit this.
\end{note}

\subsection{What the search would really decide}
\label{sec:welder-search}

The test shape is too small to show it, so take the ResNet-stage shape
(\code{[32, 64, 56, 56]}, $3\times3$, 64$\to$64). Below is the fused
conv+BN+SiLU traffic for tiles of different \emph{shapes}, with the reduction staged
8 input channels at a time (Welder's \code{rstep}, \code{rc = 8}) for the footprint:

\begin{center}
\small
\begin{tabular}{lrrrl}
\toprule
output tile $[b, c_{\text{out}}, h, w]$ & tiles & footprint & traffic & reachable by \\
\midrule
$[1, 1, 1, 16]$  & 458{,}752 & 2.0 KB  & 7{,}432 MB & \code{conv2d_forward} \\
$[1, 1, 1, 32]$  & 229{,}376 & 3.6 KB  & 6{,}549 MB & \code{conv2d_forward} \\
$[1, 16, 1, 32]$ & 14{,}336  & 7.9 KB  & 934 MB     & \code{conv2d_bn_silu_tiled} \\
$[1, 64, 1, 56]$ & 1{,}792   & 24.0 KB & 371 MB     & \code{conv2d_bn_silu_tiled} \\
$[1, 32, 8, 8]$  & 3{,}136   & 12.4 KB & 338 MB     & needs \code{BLOCK_OH} \\
$[1, 64, 4, 28]$ & 896       & 24.2 KB & \textbf{200 MB} & needs \code{BLOCK_OH} \\
\midrule
\multicolumn{4}{l}{BN+SiLU as separate kernels add, in every row} & $+102.8$ MB ($4|y|$) \\
\bottomrule
\end{tabular}
\end{center}

This table is the strongest conclusion of the report.
\begin{itemize}
\item \textbf{Tile shape decides the traffic, not fusion.} A $1\times\text{BLOCK\_OW}$
  strip re-reads all of \code{x}'s receptive field once per output channel, so traffic
  scales with $C_{\text{out}}$. Tiling output channels divides that by the channel
  block, and tiling rows shares the halo.
\item \textbf{Fusion is worth a constant $4|y|$.} That is 1.5\% of the unfused
  total at the best \code{conv2d_forward} tile, and 34\% at the best 2-D tile. The BN+SiLU fusion only
  pays off once the conv is tiled well.
\item \textbf{The DSL bounds the search.} Welder can only enumerate the blocks a
  kernel declares. \code{conv2d_forward} offers one row of this table.
  \code{conv2d_bn_silu_tiled} offers the channel-tiled rows. The best rows need a
  \code{BLOCK_OH} (\code{#[tile_loop_tile]} accepts one blocked output axis) and a
  reduction chunk (\code{#[tile_loop]} has none).
\item \textbf{\code{../Welder} reaches the bottom rows for free.} In TE every axis is
  tileable and \code{rstep} is part of the configuration. In the DSL the equivalent
  freedom must be \emph{declared}, kernel by kernel.
\end{itemize}

\subsection{Training, step by step}
\label{sec:welder-train}

Training BN normalises with \emph{batch} statistics. For a channel $c$,
$\mu_c, \sigma_c$ need every element of channel $c$ across $B, H, W$ before any
output of that channel can be produced. (Teenygrad's BN2d training path does not do
this today, see \issue{teenygrad-fq0l}. This section assumes an NCHW stats/normalize
pair like the BN1d one in \S\ref{sec:bn}.) The graph gains a node and a fan-out:
\code{conv -> {stats, normalize}}, \code{stats -> normalize},
\code{normalize -> silu}.

\textbf{Propagate.} The seed is a SiLU tile $[1,1,1,b]$:
\begin{enumerate}
\item \textbf{normalize:} with \code{N} walked, its \code{x} resolves to
  \code{[None, 1]}, the whole channel. Even a channel-tiled normalize needs its
  channel's $\mu, \sigma$, which comes from \textbf{stats}.
\item \textbf{stats:} a \code{mean} tile \code{[1]} resolves \code{x} to every
  $(b, h, w)$ of the channel. This is the reduction-domain-is-full-extent rule of
  \S\ref{sec:bn}.
\end{enumerate}
Whatever tile the output wants, \textbf{the conv tile feeding stats is a whole
channel slab} $[B, 1, OH, OW]$.

\textbf{The candidate connections}, as traffic on the intermediate (in units of
$|y|$, excluding the conv's own input reads):

{\small
\begin{xltabular}{\linewidth}{>{\raggedright\arraybackslash}p{5.2cm}c>{\raggedright\arraybackslash}X}
\toprule
\textbf{Schedule} & \textbf{$|y|$ moved} & \textbf{Verdict} \\
\midrule
\endhead
(A) four kernels & 6 & Baseline: conv writes; stats reads; normalize reads and writes; SiLU reads and writes. \\
(B) conv $|$ stats $|$ normalize+SiLU & 4 & \textbf{What Welder picks.} normalize$\to$SiLU is element-wise with matching tiles, so it fuses exactly as in inference. \\
(C) conv+partial stats $|$ normalize+SiLU & 3 + partials & Outside the model. Each conv tile would contribute partial sums to a cross-tile reduction (atomics, or a second phase). Welder's operator-tiles are independent, and there is no edge type for ``all tiles of the producer must finish first''. \\
(D) one program per channel, all four nodes & 1 & Footprint $B\cdot OH\cdot OW\cdot 4$ per program: \textbf{40 KB} on the test shape (fits 48 KB), \textbf{392 KB} on ResNet (rejected). Even where it fits, there are only $C_{\text{out}} = 8$ programs, and the wave penalty of \S4.1 rejects it on any GPU with more SMs than that. \\
\bottomrule
\end{xltabular}
}

So the training schedule in Figure~\ref{fig:schedules} cuts at both of conv's out-edges
and saves $2|y|$, against $4|y|$ in inference. The barrier between computing a
statistic and using it is a property of the math, not of the scheduler. Welder's model
(independent tiles plus connections) represents it as a mandatory global edge.

\textbf{What the DSL would have to say for (B):}
\begin{itemize}
\item an NCHW stats kernel reducing over \code{B} and \code{HW}, which needs
  \code{reduction_axis} to become a list (\S\ref{sec:bn-train});
\item a normalize kernel with a spec that resolves per-channel operands to the
  channel tile, not full extent (\code{block = 1} or the gridded-axis rule);
\item SiLU spliced into normalize's map pass. The map pass ends in a generated store
  of the body's expression, so a consumer core can be appended to that expression,
  the multi-pass analogue of the \code{finish} proposal.
\end{itemize}

\subsection{What teenygrad needs, in order}
\label{sec:needs}

\begin{enumerate}[leftmargin=1.4em,itemsep=0.35em]
\item \textbf{Grid-decoded axes are bindings with block 1} (\S\ref{sec:conv-overestimate}).
  This fixes the 91--512$\times$ conv footprint over-estimate, resolves conv's
  \code{H} window, and makes the per-channel operands of every norm resolve to one
  channel. It removes the two-spellings gotcha.
\item \textbf{A graph-level propagate with positional edges.} Reverse-topo
  \code{resolve_inputs} plus full-extent materialisation, clamping, and re-boxing of
  flattened tiles (\S\ref{sec:edges}). Rider matching in \code{GridSpec} should then
  also be positional (\issue{1nr.17}).
\item \textbf{A graph-level \code{ConstLookup}} from each node's parameters
  (\issue{1nr.30}). Without it every window resolves to full extent.
\item \textbf{Declare what is read.} Weights, per-channel parameters through raw
  pointers, and running buffers need axes. An \code{extent} that may be an expression
  (\code{C_IN / G}) would cover conv's weight.
\item \textbf{Searchable conv tiles.} A \code{BLOCK_OH} (lift the single-blocked-axis
  limit on \code{#[tile_loop_tile]}), a \code{BLOCK_N} on the generated conv like the
  tiled variant has, and a reduction chunk on \code{#[tile_loop]} (Welder's
  \code{rstep}). \S\ref{sec:welder-search} says this is worth an order of magnitude;
  fusion is worth a third.
\item \textbf{A fusion core for generated loops} = carry after \code{finish}, plus a
  broadcast-operand-aware core so BN can be spliced (\S\ref{sec:welder-infer}).
\item \textbf{Multi-axis \code{reduce}} for NCHW training BN (\S\ref{sec:bn-train}),
  then \issue{teenygrad-fq0l} and \issue{teenygrad-1nr.18.7}.
\item \textbf{The anduin pipeline itself}: \code{schedule_graph},
  \code{Trace::trace_graph}, a profiler, and codegen, which are all \code{todo!()}
  today.
\end{enumerate}
